\chapter{Symbolic Reduction and Equivalence Testing\label{chap:SymbolicEquivalenceTesting}}

In this chapter, we define two forms of symbolic reasoning ---
\emph{symbolic reduction} and \emph{symbolic equivalence testing}.
Symbolic reduction simplifies expressions into \equivalentexprsterm{}
that are simpler to reason about.
In our context, equivalence means that we can substitute one expression for another without
affecting the semantics of the overall specification.
%
Symbolic equivalence is a \emph{conservative} test.
By conservative, we mean that if a test for equivalence returns $\True$ then the expressions
being compared are indeed equivalent, but if the test returns $\False$ then
there are two possibilities:
\begin{itemize}
  \item the expressions are not equivalent;
  \item the expressions are equivalent, but the reasoning power of our rules
  is not enough to prove it, and so we conservatively answer negatively.
\end{itemize}
In proof-theoretic terms, we say that our equivalence test is \emph{sound} but \emph{incomplete}.

Notice that for a conservative test, it is always correct to return $\False$.

In this chapter, we use the special constant value \RenderConstant{CannotBeTransformed}
to represent a failure in transforming an expression into a desired form (the specific
desired form varies according to the functions utilizing this value).

\ChapterOutline
\begin{itemize}
  \item \secref{Equivalence Definitions} formally defines equivalence for expressions, types, constraints, and bitfields;
  \item \secref{symbolicexpressions} defines \symbolicexpressionsterm{} and operations over \symbolicexpressionsterm{}; and
  \item \secref{SymbolicReductionAndEquivalenceTestingRules} defines rules for symbolic reduction and equivalence testing.
\end{itemize}

\section{Equivalence Definitions\label{sec:Equivalence Definitions}}
\begin{definition}[Expression Equivalence]
\hypertarget{def-equivalentexprsterm}{}
We say that $\veone$ and $\vetwo$ are \emph{\equivalentexprsterm{}} in the \staticenvironmentterm{} $\tenv$
if for every dynamic environment $\denv$
and environment $\env = (\tenv, \denv)$, the following condition holds:
\[
\{ C \;|\; \evalexpr(\env, \veone) \evalarrow C \} = \{ C \;|\; \evalexpr(\env, \vetwo) \evalarrow C \} \enspace.
\]
\end{definition}

\begin{definition}[Type Equivalence]
\hypertarget{def-equivalenttypesterm}{}
We say that $\vt$ and $\vs$ are \emph{\equivalenttypesterm{}} in the \staticenvironmentterm{} $\tenv$
if both types have the same AST label,
and for every dynamic environment $\denv$
and environment $\env = (\tenv, \denv)$, the following condition holds:
\[
\dynamicdomain(\env, \vt) = \dynamicdomain(\env, \vs) \enspace.
\]
\end{definition}

\begin{definition}[Constraint Equivalence]
\hypertarget{def-equivalentconstraintsterm}{}
We say that $\vcone$ and $\vctwo$ are \emph{\equivalentconstraintsterm{}} in the \staticenvironmentterm{} $\tenv$
if for every dynamic environment $\denv$
and environment $\env = (\tenv, \denv)$, the following condition holds:
\[
\dynamicdomain(\env, \vcone) = \dynamicdomain(\env, \vctwo) \enspace.
\]
\end{definition}

Bitfield equivalence is defined by \TypingRuleRef{BitFieldEqual} where
the conservative tests for expression equivalence ($\exprequal$)
and type equivalence ($\typeequal$) are replaced by their respective
definitions above.

\section{Symbolic Expressions\label{sec:symbolicexpressions}}
\hypertarget{def-symbolicexpressionterm}{}
Our symbolic reduction and equivalence testing rules employ \emph{\symbolicexpressionsterm}, defined below:
\RenderTypes{symbolic_expressions}

These \symbolicexpressionsterm{} provide a canonical representation for multivariate polynomials
(that is, polynomials over multiple variables). The motivation for choosing this representation
is examples similar to the one in \listingref{IR}, which require assigning the
multivariate polynomials \verb|8*M*N| and \verb|M*N*8| a unique representation.
\ASLListing{Example requiring a canonical polynomial representation}{IR}{\definitiontests/IR.asl}

We now explain each component of a \symbolicexpressionterm{} and how it can be interpreted as a mathematical formula
via the interpretation function $\alpha$.
We also define operations over \symbolicexpressionsterm{}.

\begin{definition}[Unitary Monomial]
A \emph{Unitary Monomial} is a partial function from identifiers to positive integers\footnote{A unitary monomial has a unit factor,
for example $x^3$, whereas a non-unitary monomial has a non-unit factor, for example, $2 x^3$.}.

A non-empty unitary monomial, $\vm\in\unitarymonomial$ where $\vm \neq \emptyfunc$, can be interpreted as follows:
\[
  \alpha(\vm) \triangleq \prod_{\vx \in \dom(\vm)} \vx^{\vm(\vx)} \enspace.
\]

An empty unitary monomial is interpreted as the constant $1$:
\[
  \alpha(\emptyfunc) \triangleq 1 \enspace.
\]
\end{definition}
For example,
\[
  \alpha(\ \{\vx\mapsto 3, \vy\mapsto 1, \vz\mapsto2\}\ ) = x^3 \cdot y \cdot z^2 \enspace.
\]

\RenderRelation{mul_monomials}

\begin{mathpar}
  \inferrule{
    {
      \vf \eqdef \lambda \vx\in\Identifier.\
      \left\{
      \begin{array}{ll}
        \vfone(\vx) & \text{if } \vx \in \dom(\vfone) \setminus \dom(\vftwo)\\
        \vftwo(\vx) & \text{if } \vx \in \dom(\vftwo) \setminus \dom(\vfone)\\
        \vfone(\vx)+\vftwo(\vx) & \text{if } \vx \in \dom(\vfone) \cap \dom(\vftwo)\\
      \end{array}
      \right.
    }
  }
  {
    \mulmonomials(\overname{\vfone}{\vmone}, \overname{\vftwo}{\vmtwo}) \typearrow \overname{\vf}{\vm}
  }
\end{mathpar}
For example,
\[
\mulmonomials(\{\vx\mapsto{}3, \vy\mapsto{}1, \vz\mapsto{}2\}, \{\vx\mapsto{}1, \vw\mapsto{}2\}) =
              \{\vx\mapsto{}4, \vy\mapsto{}1, \vz\mapsto{}2, \vw\mapsto{}2\}
\]

\begin{definition}[Polynomial]
  \emph{Polynomials} are partial functions from unitary monomials to rationals other than zero.
  Intuitively, each unitary monomial is mapped to its factor in the polynomial.
  A polynomial $\vp$ can be interpreted as follows:
  %
\[
  \alpha(\vp) \triangleq \sum_{\vm \in \dom(\vp)} \vp(\vm)\cdot\alpha(\vm)
\]
\end{definition}
For example,
\[
  \left(\left\{
    \begin{array}{lcl}
      \{\vx\mapsto 3, \vy\mapsto 1, \vz\mapsto2\} &\mapsto& -1,\\
      \{\vx\mapsto 2, \vy\mapsto 1\} &\mapsto& \frac{3}{4}
    \end{array} \right\}\right) =
    -1\cdot x^3 \cdot y \cdot z^2 + \frac{3}{4} \cdot \vx^2\cdot \vy \enspace.
\]

\RenderRelation{polynomial_of_int}
It is defined as follows:
\[
\polynomialofint(\vi) \typearrow \overname{\{\emptyfunc \mapsto \vi\}}{\vp} \enspace.
\]

\TypingRuleDef{PolynomialOfVar}
\RenderRelation{polynomial_of_var}
% NO_EXAMPLE
\RenderProseAndFormally{polynomial_of_var}

\TypingRuleDef{AddPolynomials}
\RenderRelation{add_polynomials}

See \ExampleRef{Adding a List of Polynomials}

\begin{mathpar}
\inferrule{
  {
    \begin{array}{l}
    \vf \eqdef \lambda \vm\in\unitarymonomial.\\
    \left\{
    \begin{array}{ll}
      \vfone(\vm) & \text{if } \vm \in \dom(\vfone) \setminus \dom(\vftwo)\\
      \vftwo(\vm) & \text{if } \vm \in \dom(\vftwo) \setminus \dom(\vfone)\\
      \bot        & \text{if } \vm \in \dom(\vfone) \cap \dom(\vftwo) \text{ and } \vfone(\vm)+\vftwo(\vm) = 0\\
      \vfone(\vm)+\vftwo(\vm) & \text{if } \vm \in \dom(\vfone) \cap \dom(\vftwo) \text{ and } \vfone(\vm)+\vftwo(\vm) \neq 0\\
    \end{array}
    \right.
    \end{array}
  }
}{
  \addpolynomials(\overname{\vfone}{\vpone}, \overname{\vftwo}{\vptwo}) \typearrow \overname{\vf}{\vp}
}
\end{mathpar}

\TypingRuleDef{AddPolynomialList}
The overloaded function
\[
  \addpolynomials : \KleeneStar{\polynomial} \rightarrow \polynomial
\]
adds a list of polynomials.

\ExampleDef{Adding a List of Polynomials}
For example, if $\alpha(\vpone) = 3\cdot{}x^2\cdot{}y$,
$\alpha(\vptwo) = 4\cdot{}x^2\cdot{}y$, and
$\alpha(\vpthree) = x\cdot{}y^2$,
then $\alpha(\addpolynomials([\vpone, \vptwo, \vpthree])) = 7\cdot{}x^2\cdot{}y + x\cdot{}y^2$.

\begin{mathpar}
\inferrule[empty]{}{ \addpolynomials(\emptylist) \typearrow \emptyfunc }
\hva\and
\inferrule[one]{}{ \addpolynomials([ \vp ]) \typearrow \vp }
\hva\and
\inferrule[two\_or\_more]{
  \addpolynomials(\vp_{2..k}) \typearrow \vpp\\
  \addpolynomials(\vp_1, \vpp) \typearrow \vp
}{
  \addpolynomials(\vp_{1..k}) \typearrow \vp
}
\end{mathpar}

\RenderRelation{mul_polynomials}

\begin{mathpar}
\inferrule{
  {
    \vps \eqdef \left\{
      \begin{array}{l}
        \{\mulmonomials(\vmone, \vmtwo) \mapsto \vfone(\vmone)\times\vftwo(\vmtwo)\}  \;|\; \\
       \vmone\in\dom(\vfone)\ \land\ \vmtwo\in\dom(\vftwo)
    \end{array}
    \right\}
  }\\
  \orderedps \eqdef \listset{\vps}\\
  \addpolynomials(i=1..k: \orderedps) \typearrow \vp\\
}{
  \mulpolynomials(\overname{\vfone}{\vpone}, \overname{\vftwo}{\vptwo}) \typearrow \vp
}
\end{mathpar}

\RenderRelation{polynomial_divide_by_term}

\begin{mathpar}
\inferrule[divides]{
  \vptwo \eqdef [\vm \mapsto \vf]\\
  \vpone = \mulpolynomials(\vp, \vptwo)
}{
  \polynomialdividebyterm(\vpone, \vm, \vf) \typearrow \vp
}
\end{mathpar}

\begin{mathpar}
\inferrule[does\_not\_divide]{
  \vptwo \eqdef [\vm \mapsto \vf]\\
  \forall \vp\in\polynomial.\ \vpone \neq \mulpolynomials(\vp, \vptwo)
}{
  \polynomialdividebyterm(\vpone, \vm, \vf) \typearrow \CannotBeTransformed
}
\end{mathpar}

For example,
if $\alpha(\vpone) = 6\cdot{}x^2\cdot{}y + 2\cdot{}x\cdot{}z$
and $\alpha(\vm) = x$ and $\vf=2$ then
\[
\alpha(\polynomialdividebyterm(\vpone, \vm, \vf)) =  3\cdot{}x\cdot{}y + z \enspace.
\]
%
In contrast, if $\alpha(\vpone) = x$
and $\alpha(\vm) = x^2$ and $\vf=1$ then
\[
\alpha(\polynomialdividebyterm(\vpone, \vm, \vf)) =  \CannotBeTransformed \enspace.
\]

Notice that for all $\vpone$ and $\vm$ the following holds:
\[
\alpha(\polynomialdividebyterm(\vpone, \vm, 0)) =  \CannotBeTransformed \enspace.
\]

\section{Typing Rules\label{sec:SymbolicReductionAndEquivalenceTestingRules}}
We employ the following rules:
\begin{itemize}
  \item \TypingRuleRef{Normalize}
  \item \TypingRuleRef{ReduceConstraint}
  \item \TypingRuleRef{ReduceConstraints}
  \item \TypingRuleRef{ToIR}
  \item \TypingRuleRef{ExprEqual}
  \item \TypingRuleRef{ExprEqualNorm}
  \item \TypingRuleRef{ExprEqualCase}
  \item \TypingRuleRef{TypeEqual}
  \item \TypingRuleRef{BitwidthEqual}
  \item \TypingRuleRef{BitFieldsEqual}
  \item \TypingRuleRef{BitFieldEqual}
  \item \TypingRuleRef{ConstraintEqual}
  \item \TypingRuleRef{SlicesEqual}
  \item \TypingRuleRef{SliceEqual}
  \item \TypingRuleRef{ArrayLengthEqual}
\end{itemize}

\TypingRuleDef{Normalize}
\RenderRelation{normalize}

\hypertarget{def-normalizedterm}{}
We refer to an expression in the image of $\normalize$ as a \emph{\normalizedexpressionterm}.

\ExampleDef{Normalize}
The specification in \listingref{Normalize} shows examples of expressions (on the right-hand-sides of assignments and declarations)
and their corresponding normalized expressions.
\ASLListing{Normalizing expressions}{Normalize}{\typingtests/TypingRule.Normalize.asl}


\RenderProseAndFormally{normalize}
\CodeSubsection{\NormalizeBegin}{\NormalizeEnd}{../StaticModel.ml}

\TypingRuleDef{ReduceConstraint}
\RenderRelation{reduce_constraint}

\ExampleDef{Symbolically Simplifying Constraints}
The specification in \listingref{ReduceConstraint} shows examples of constraints
and their simplified counterparts.
\ASLListing{Symbolically simplifying constraints}{ReduceConstraint}{\typingtests/TypingRule.ReduceConstraint.asl}


\RenderProseAndFormally{reduce_constraint}
\CodeSubsection{\ReduceConstraintBegin}{\ReduceConstraintEnd}{../StaticOperations.ml}

\TypingRuleDef{ReduceConstraints}
\RenderRelation{reduce_constraints}

See \ExampleRef{Symbolically Simplifying Constraints}.


\RenderProseAndFormally{reduce_constraints}
\CodeSubsection{\ReduceConstraintsBegin}{\ReduceConstraintsEnd}{../StaticOperations.ml}

\TypingRuleDef{ToIR}
\RenderRelation{to_ir}

\ExampleDef{Transforming Expressions into Symbolic Expressions}
The specification in \listingref{ToIR} shows examples of expressions (on the right-hand sides of assignments and declarations)
and their corresponding symbolic representations (in comments, as ASL expressions).
\ASLListing{Transforming expressions into symbolic expressions}{ToIR}{\typingtests/TypingRule.ToIR.asl}


We make use of the following helper relations:
\RenderRelation{is_div_binop}
It defined as follows:
\RenderProseAndFormally{is_div_binop}
\RenderRelation{scale_polynomial}
It is defined as follows:
\RenderProseAndFormally{scale_polynomial}

\RenderProseAndFormally{to_ir}
\CodeSubsection{\ToIRBegin}{\ToIREnd}{../StaticModel.ml}

\TypingRuleDef{ExprEqual}
\RenderRelation{expr_equal}

\ExampleDef{Expression Equivalence}
The specification in \listingref{ExprEqual} is well-typed as the expressions
comprising the array lengths are equivalent, as detailed in comments.
\ASLListing{Equivalent expressions}{ExprEqual}{\typingtests/TypingRule.ExprEqual.asl}


\RenderProseAndFormally{expr_equal}

\TypingRuleDef{ExprEqualNorm}
\RenderRelation{expr_equal_norm}

See \ExampleRef{Expression Equivalence}.


\RenderProseAndFormally{expr_equal_norm}

\TypingRuleDef{ExprEqualCase}
\RenderRelation{expr_equal_case}

See \ExampleRef{Expression Equivalence}.
In the following definition recall that a conjunction over an empty set equals $\True$.

\RenderProseAndFormally{expr_equal_case}

\TypingRuleDef{TypeEqual}
\RenderRelation{type_equal}

\ExampleDef{Equivalent Types}
The specification in \listingref{TypeEqual} is well-typed.
Specifically, the types of the arguments for the \verb|impdef| declaration
of \verb|Foo| match types of the corresponding arguments for the overridden declaration
of \verb|Foo|.
\ASLListing{Equivalent Types}{TypeEqual}{\typingtests/TypingRule.TypeEqual.asl}


\RenderProseAndFormally{type_equal}

\TypingRuleDef{PatternEqual}
\RenderRelation{pattern_equal}

See \ExampleRef{Expression Equivalence}.

\RenderProseAndFormally{pattern_equal}

\TypingRuleDef{BitwidthEqual}
\RenderRelation{bitwidth_equal}

\ExampleDef{Equivalence of Bitwidth Expressions}
\listingref{BitwidthEqual} shows examples of equivalent bitwidth expressions on both sides
of the respective local declaration statements.
\ASLListing{Equivalence of bitwidth expressions}{BitwidthEqual}{\typingtests/TypingRule.BitwidthEqual.asl}

\RenderProseAndFormally{bitwidth_equal}

\TypingRuleDef{BitFieldsEqual}
\RenderRelation{bitfields_equal}

\ExampleDef{Equivalent Lists of Bitfields}

The specification in \listingref{BitFieldsEqual} is well-typed.
Specifically, the lists of bitfields for the \verb|bv| argument in both signatures
of \verb|Foo| are equivalent.
\ASLListing{Equivalent lists of bitfields}{BitFieldsEqual}{\typingtests/TypingRule.BitFieldsEqual.asl}

The specification in \listingref{BitFieldsEqual-bad} is ill-typed,
the lists of bitfields for the \verb|bv| argument in both signatures
of \verb|Foo| are not equivalent.
\ASLListing{Non-equivalent lists of bitfields}{BitFieldsEqual-bad}{\typingtests/TypingRule.BitFieldsEqual.bad.asl}


\RenderProseAndFormally{bitfields_equal}

\TypingRuleDef{BitFieldEqual}
\RenderRelation{bitfield_equal}

\ExampleDef{Bitfield Equivalence}
In \listingref{BitFieldEqual}, the bitfields of the \bitvectortypesterm{} on the left-hand-side
of assignments and the bitfields of the same name in the \bitvectortypesterm{} on the right-hand-side
of the same assignment are equivalent.
\ASLListing{Equivalent bitfields}{BitFieldEqual}{\typingtests/TypingRule.BitFieldEqual.asl}

The specifications in \listingref{BitFieldEqual-bad1} and \listingref{BitFieldEqual-bad2}
show cases where the \verb|data| bitfields of the \bitvectortypesterm{} on the two sides
of an assignment are not equivalent. This is due to either having different slices
or different nested bitfields.
\ASLListing{Bitfields with differing slices}{BitFieldEqual-bad1}{\typingtests/TypingRule.BitFieldEqual.bad1.asl}
\ASLListing{Bitfields with differing nested bitfields}{BitFieldEqual-bad2}{\typingtests/TypingRule.BitFieldEqual.bad2.asl}


\RenderProseAndFormally{bitfield_equal}

\TypingRuleDef{ConstraintsEqual}
\RenderRelation{constraints_equal}

\ExampleDef{Equivalent Lists of Constraints}
The specification in \listingref{ConstraintsEqual}
shows two lists of constraints --- the ones used to type \verb|x| and the ones used to type \verb|y|.
The expression \verb|x as integer{w..2 * w + z, w + w + w}|
compares the list \verb|3 * w, w..w + z + w| to the list \verb|w..2 * w + z, w + w + w|.
Since the corresponding constraints are determined to be equivalent,
the expression is well-typed.
The expression \verb|x as integer{w + w + w, w..2 * w + z}| on the other hand
has the constraints in a different order, which means that the two lists of constraints
are not considered equivalent. This is still considered well-typed (see \TypingRuleRef{CheckATC}),
and the type assertion will be checked dynamically (see \SemanticsRuleRef{ATC}).
\ASLListing{Equivalent lists of constraints}{ConstraintsEqual}{\typingtests/TypingRule.ConstraintsEqual.asl}


\RenderProseAndFormally{constraints_equal}

\TypingRuleDef{ConstraintEqual}
\RenderRelation{constraint_equal}

\ExampleDef{Equivalent Constraints}
The specification in \listingref{ConstraintEqual} is well-typed,
since each constraint on a right-hand-side of an assignment is equivalent to the
corresponding constraint on the left-hand-side of that assignment.
\ASLListing{Equivalent constraints}{ConstraintEqual}{\typingtests/TypingRule.ConstraintEqual.asl}


\RenderProseAndFormally{constraint_equal}

\TypingRuleDef{SlicesEqual}
\RenderRelation{slices_equal}

\ExampleDef{Equivalent Lists of Slices}
\listingref{SlicesEqual} shows an example of two equivalent lists of slices.
\ASLListing{Equivalent lists of slices}{SlicesEqual}{\typingtests/TypingRule.SlicesEqual.asl}


\RenderProseAndFormally{slices_equal}

\TypingRuleDef{SliceEqual}
\RenderRelation{slice_equal}

This function operates both on \typedast{} nodes, where both slices are only
of the form $\SliceLength(\ \cdot\ ,\ \cdot\ )$, as well as \untypedast{} nodes where all slice forms
are possible. The latter is used during checking overriding subprograms (\secref{Overriding Subprograms}).

\ExampleDef{Equivalent Slices}
The specification in \listingref{SliceEqual} shows examples of equivalent slices.
\ASLListing{Equivalent slices}{SliceEqual}{\typingtests/TypingRule.SliceEqual.asl}

In \listingref{overriding}, the slice for the bitfield \verb|lsb| in both the \texttt{impdef} version
of \verb|Foo| and the overridden version of \verb|Foo| are equivalent.


\RenderProseAndFormally{slice_equal}

\TypingRuleDef{ArrayLengthEqual}
\RenderRelation{array_length_equal}

\ExampleDef{Array Length Expression Equivalence}
See \ExampleRef{Expression Equivalence} for examples of equivalent array length expressions.

The specification in \listingref{ArrayLengthEqual-bad} is ill-typed since the array length
expressions are of different kinds (\integertypeterm{} and \enumerationtypeterm),
even though they both contain three elements each.
\ASLListing{Non-equivalent array length expressions}{ArrayLengthEqual-bad}{\typingtests/TypingRule.ArrayLengthEqual.bad.asl}


\RenderProseAndFormally{array_length_equal}

\TypingRuleDef{PolynomialToExpr}
\RenderRelation{polynomial_to_expr}

See \ExampleRef{Transforming Expressions into Symbolic Expressions}.

\texthypertarget{relation-sortmonomialbindings}
The helper function $\sortmonomialbindings{f}$ sorts the list of monomial-to-factor bindings in $f$
using the relation $\comparemonomialbindings$.

\RenderProseAndFormally{polynomial_to_expr}

\TypingRuleDef{CompareMonomialBindings}
\RenderRelation{compare_monomial_bindings}

\ExampleDef{Comparing Monomial Bindings}
Comparing $(\vx, -1)$ to $(\vx, -1)$ yields $0$.

Comparing $(\vx, 1)$ to $(\vy, 1)$ yields $-1$ as $\vx$ comes before $\vy$ and $\vx$ is the
first identifier for which the two monomials differ.

The function \RenderConstant{compare_identifier} compares two identifiers, which are lists of ASCII characters,
via the lexicographic ordering.

\RenderProseAndFormally{compare_monomial_bindings}

\TypingRuleDef{MonomialsToExpr}
\RenderRelation{monomials_to_expr}

\ExampleDef{Transforming a List of Unitary Monomials}
Transforming the list of unitary monomials whose formulae are given by
$(\vx \times \vy^2, 2)$,
$(\vx \times \vy^2, 3)$ and $(\vx \times \vy, -5)$
yields the expression \verb|5 * x * y * y - 5 * x * y|.


\RenderProseAndFormally{monomials_to_expr}

\TypingRuleDef{MonomialToExpr}
\RenderRelation{monomial_to_expr}

\ExampleDef{Transforming Expressions and Rationals}
The following are examples of inputs and outputs of $\monomialtoexpr$:
\[
\monomialtoexpr(\EVar(\vx), 0) \typearrow (\ELiteral(\LInt(0)), 0)
\]

\[
\monomialtoexpr(\EVar(\vx), 5) \typearrow (\AbbrevEBinop{*}{\ELInt{5}}{\AbbrevEVar{\vx}}, 1)
\]

\[
\monomialtoexpr(\EVar(\vx), 9/12) \typearrow
(\AbbrevEBinop{/}
{ (\AbbrevEBinop{*}{\AbbrevEVar{\vx}}{3}) }
{ \ELInt{4} }, 1)
\]

\[
\monomialtoexpr(\EVar(\vx), -1) \typearrow (\EUnop(\NEG, \AbbrevEVar{\vx}), -1)
\]



\RenderProseAndFormally{monomial_to_expr}

\TypingRuleDef{SymAddExpr}
\RenderRelation{sym_add_expr}

The effect of the function can be summarized by the following table:
\begin{center}
\begin{tabular}{|c|c|c|c|}
\hline
& \multicolumn{3}{c|}{$\vsone$}\\
\hline
$\vstwo$                      & -1                         & 0               & 1\\
\hline
-1                            &  $(\veone+\vetwo, \vsone)$ & (\vetwo,\vstwo) & $(\veone-\vetwo, \vsone)$ \\
\hline
 0                            &  $(\veone, \vsone)$        & (\veone,\vsone) & $(\veone, \vsone)$ \\
 \hline
 1                            &  $(\veone-\vetwo, \vsone)$ & (\vetwo,\vstwo) & $(\veone+\vetwo, \vsone)$ \\
\hline
\end{tabular}
\end{center}

\ExampleDef{Symbolically Summing Signed Expressions}
The following are some examples of inputs and outputs for symbolic addition:
\begin{center}
\begin{tabular}{rrrrrr}
\textbf{$\vsone$} & \textbf{$\veone$} & \textbf{$\vstwo$} & \textbf{$\vetwo$} & \textbf{$\ve$} & \textbf{$\vs$}\\
\hline
-1 & \verb|x| & -1 & \verb|y| & \verb|x+y| & -1\\
-1 & \verb|x| & 0 & \verb|y| & \verb|x| & -1\\
-1 & \verb|x| & 1 & \verb|y| & \verb|x-y| & -1\\
1 & \verb|x| & -1 & \verb|y| & \verb|x-y| & 1\\
\end{tabular}
\end{center}


\RenderProseAndFormally{sym_add_expr}

\TypingRuleDef{UnitaryMonomialsToExpr}
\RenderRelation{unitary_monomials_to_expr}

\ExampleDef{Transforming Unitary Monomials to Expressions}
The specification in \listingref{UnitaryMonomialsToExpr}
shows examples of unitary monomials (as right-hand-side of assignments)
and the corresponding expressions, in comments.
\ASLListing{Transforming unitary monomials to expressions}{UnitaryMonomialsToExpr}{\typingtests/TypingRule.UnitaryMonomialsToExpr.asl}


\RenderProseAndFormally{unitary_monomials_to_expr}

\TypingRuleDef{SymMulExpr}
\RenderRelation{sym_mul_expr}

\ExampleDef{Symbolic Multiplication}
The following table shows examples of symbolically multiplying the expression $\veone$ by the expression $\vetwo$,
yielding the result in $\ve$:
\begin{center}
\begin{tabular}{lll}
\textbf{$\veone$} & \textbf{$\vetwo$} & \textbf{$\ve$}\\
\hline
\verb|x*0|  & \verb|1|    & \verb|x*0|\\
\verb|1|    & \verb|x*0|  & \verb|x*0|\\
\verb|x*0|  & \verb|x*0|  & \verb|(x*0)*(x*0)|
\end{tabular}
\end{center}


\RenderProseAndFormally{sym_mul_expr}

\TypingRuleDef{TypeOf}
\RenderRelation{type_of}

\ExampleDef{The Type Bound to an Identifier}
In \listingref{expressions-evar}, the type bound to the global storage element
\verb|global_non_constant| after annotating the global storage element declarations
is \verb|integer{19}|,
%
the type bound to the local storage element \verb|var_x| in the \staticenvironmentterm{}
after annotating the local declaration \verb|var var_x = LOCAL_CONSTANT;| is \verb|integer{7}|.
%
Invoking $\typeof$ for the \staticenvironmentterm{} just before annotating the local
declaration \verb|var x = t;| in \listingref{expressions-evar-undefined} yields a \typingerrorterm.


\RenderProseAndFormally{type_of}
\CodeSubsection{\TypeOfBegin}{\TypeOfEnd}{../StaticModel.ml}

\TypingRuleDef{NormalizeOpt}
\RenderRelation{normalize_opt}

See \ExampleRef{Normalize}.


\RenderProseAndFormally{normalize_opt}
